\section{\ours Environment}
\label{sec:environment}
\vspace{-5pt}

In this section, we will introduce the task definition of autonomous agents, the components and implementation of the \ours environment, and the supported observation and action spaces.

\subsection{Task Definition}
An autonomous digital agent task can be formalized as a partially observable Markov decision process (POMDP) $(\mathcal{S}, \mathcal{O}, \mathcal{A}, \mathcal{T}, \mathcal{R})$ 
with state space $\mathcal{S}$, observation space $\mathcal{O}$ (\S\ref{observation_space}, including natural language $\mathcal{I}$), action space $\mathcal{A}$ (\S\ref{action_space}), transition function $\mathcal{T}: \mathcal{S} \times \mathcal{A} \to \mathcal{S}$, and reward function $\mathcal{R}: \mathcal{S} \times \mathcal{A} \to \mathbb{R}$.
Given current observation $o_{t} \in \mathcal{O}$ (a natural language instruction observation and a screenshot observation (\textit{e.g.}, computer screenshot), accessibility (a11y) tree, or their combination according to facilities available), an agent generates executable action $a_t \in \mathcal{A}$ (\textit{e.g.}, clicking on the certain pixel of the screen --- \texttt{.click(300, 540, button=`right')}, press key combination --- \texttt{.hotkey(`ctrl', `alt', `t')}), which results in a new state $s_{t+1} \in \mathcal{S}$ (\textit{e.g.}, current Desktop environment) and a new partial observation $o_{t+1} \in \mathcal{O}$ (\textit{e.g.}, current screenshot).
The interaction loop repeats until an action that marks termination (\verb|DONE| or \verb|FAIL|, see Sec.~\ref{action_space}) is generated or the agent reaches the max number of steps (\textit{e.g.}, 15 in our experiments).
In this version of \ours, we implement an execution-based reward function $\mathcal{R}: \mathcal{S} \times \mathcal{A} \to [0, 1]$ (\S\ref{eval}). 
The reward function awards a value of 1 or a positive decimal under 1 at the final step if the state transitions meet the expectations of the task objective (i.e., the goal is successfully achieved or partially achieved), or if the agent accurately predicts failure for an infeasible task.
In all other scenarios, it returns 0.

\subsection{Real Computer Environment Infrastructure}
\vspace{-5pt}

\ours is an executable and controllable environment that supports task initialization, execution-based evaluation, and interactive agent learning in a range of \textit{real} operating systems (\textit{e.g.},~Ubuntu, Windows, macOS) using virtual machine techniques, shown in the middle and right of Fig.~\ref{fig:arch_figure}. 
Virtual machine offers a safe isolated environment and prevents the agent resulting in irreversible damaging effect on the real host machine. 
The snapshot feature also enables efficient reset of the virtual environment.
The environment is configured through a config file (shown in the left of Fig.~\ref{fig:arch_figure}) for interface initialization during the initialization phase (including downloading files, opening software, adjusting interface layout) (\S\ref{init_setup}, highlighted with red in Fig.~\ref{fig:arch_figure}), post-processing during the evaluation phase (activating certain windows, saving some files for easy retrieval of information, highlighted with orange), and acquiring files and information for evaluation (such as the final spreadsheet file for spreadsheet tasks, cookies for Chrome tasks, highlighted with yellow in Fig.~\ref{fig:arch_figure}), as well as the evaluation functions and parameters used (\S\ref{eval}, highlighted with green in Fig.~\ref{fig:arch_figure}).
See App.~\ref{app:env_infra} for more details.

\begin{figure}[htbp]
    \centering
    \includegraphics[width=\linewidth]{images/arch_figure_w_config_new.pdf}
    \caption{
    Overview of the \ours environment infrastructure. The environment uses a configuration file for initializing tasks (highlighted in red), agent interaction, post-processing upon agent completion (highlighted in orange), retrieving files and information (highlighted in yellow), and executing the evaluation function (highlighted in green). 
    Environments can run in parallel on a single host machine for learning or evaluation purposes. 
    Headless operation is supported.
    }
    \vspace{-15pt}
    \label{fig:arch_figure}
\end{figure}

\subsubsection{Overview}
\vspace{-5pt}

\ours environment runs on the host machine.
Its Coordinator accepts a configuration file at the initialization of a computer task, runs commands to automatically create a virtual machine instance, and initializes the required state for the task through the Task Manager.
The configuration file specifies the snapshot of the virtual machine to be used (which stores the complete state of a computer at a certain moment and can be restored to this state at any time) and also indicates the information needed for setup (such as downloading files and opening some software, making some additional settings, etc.).
Once the environment is set up, agents start to interact with the environment, receiving observations such as screenshots, the accessibility (a11y) tree, and customized streams such as terminal outputs.
Agents subsequently generate executable actions (\textit{e.g.},~\texttt{.click(300, 540)}) that manipulate the keyboard and mouse.
Each action of the agent is input into the environment as a code string, and the environment's Simulator executes them in the virtual machine. 
After the completion of a task, the Task Manager performs post-processing (such as file saving, or reopening certain apps) according to the task's post-config, retrieves data to the host machine (fetching images or configuration files from the virtual machine or cloud, etc.), and then runs evaluation scripts to assess the completion of the task.
Multiple virtual machines can run simultaneously on a single host machine, thereby parallelizing training and evaluation.

\subsubsection{Initial Task Environment Setup}
\label{init_setup}
Many real-world scenarios requiring assistance occur not at the beginning of digital activities, such as right after launching an application or when a computer has just been started, but rather at intermediate stages, such as when certain software is already open or the computer has experienced a crash. 
Therefore, we aim to simulate these intermediate states as closely as possible to replicate real-world scenarios.
The naturalness we bring in also leads to more challenges for agents to model and explore.
We adopted a hybrid approach for configuration instead of solely relying on example-wise snapshots for restoration since it would store much unnecessary hardware state information, resulting in each example requiring gigabytes of space.
The procedure is divided into three stages: start the VM emulator, prepare files (download the files or scripts from the cloud, \textit{etc.} optional), and execute reprocessing commands (open files or tabs, change the window size, \textit{etc.} optional).
We provide convenient APIs to configure initial conditions and world settings, standardizing our tasks to make this process user-friendly and easily extendable for scaling.
For more details on setup see App.~\ref{app:initial_state_setup_details}.

\subsubsection{Execution-Based Evaluation}
\label{eval}

Evaluating the successful execution of general computer tasks presents a significant challenge, as these tasks defy reduction to a uniform pattern or measurement by a single metric. 
To ensure a thorough assessment, we design example-specific evaluation metrics including pre-setup, post-processing, and dedicated functions, tailored to the software in use and the task's specific requirements.
This involves interpreting the software's internal files, utilizing specific packages, and preemptively setting up scaffolding based on the software's permissions (\textit{e.g.}, opening remote debugging ports for Chrome and VLC, creating extensions for VS Code). 
Occasionally, this process may also require assistance from reverse engineering tools, such as for decrypting account information in Thunderbird.

As a result, we construct a vast collection of functions that make final wrangling and retrieve files and data information of varying types, categories, and granularities from the cloud and software from virtual machines as well as evaluation functions covering different aspects and their combinations, inputting this information as parameters to assess the outcomes.
We show some evaluation examples in Tab.~\ref{tab:evaluation_examples}.
, demonstrate the retrieval of cookie data from virtual machines, obtaining files from both virtual machines and cloud services, fetching the current runtime interface's accessibility tree from the virtual machines, and determining success based on this information whether Amazon's cookies have been deleted, whether the generated table is accurate, and whether the correct interface has been accessed.
Need to note when the type of task has real-time characteristics (such as the number of citations of someone's paper, the content of blogs, \textit{etc}.), we include dynamic functions (such as crawler scripts) inside getter to obtain the real-time values at the moment of evaluation and then use them to compare with the results obtained by the agent upon task completion.
See more in App.~\ref{app:evaluation_configuration_details}.

\begin{table}[t]
\vspace{-10pt}
\caption{
Examples of our annotated evaluation scripts, which involve retrieving data from configuration files, the environment, and the cloud, and executing functions to assess functional correctness and obtain results.
The example-wise evaluation facilitates the diversity of tasks and reliable evaluation of complex, real-world, open-ended tasks.
}
\resizebox{\textwidth}{!}{
\begin{tabular}{lp{.35\textwidth}l}
  \toprule
  \textbf{Initial State} &   \textbf{Task Instruction}  & \textbf{Evaluation Script (Simplified)}  \\
  \midrule
  \multirow{3}{*}{\raisebox{-1cm}{\includegraphics[width=5cm]{images/cases/chrome_clean_cookies.png}}}  & \multirow{3}{*}{\begin{minipage}{.35\textwidth}
  \textit{Can you help me clean up my computer by getting rid of all the cookies that Amazon might have saved?}
  \end{minipage}} &  \texttt{cookie\_data = get\_cookie\_data(env)}\\
    &   &  \texttt{rule = \{"type":"domains", }\\
    &   &  \texttt{"domains":[".amazon.com"]\}} \\
    &   &  \texttt{is\_cookie\_deleted(cookie\_data, rule)} \\ 
    &   &   \\
    &   &   \\
    &   &   \\
    &   &   \\
  \midrule
   \multirow{4}{*}{\raisebox{-1cm}{\includegraphics[width=5cm]{images/evaluation_scripts_showcase/calc_rename.png}}} &  
   \multirow{4}{*}{
   \begin{minipage}{.35\textwidth}
   \textit{Rename ``Sheet 1'' to ``LARS Resources''. Then make a copy of it. Place the copy before ``Sheet 2'' and rename it by appending a suffix ``(Backup)'', ...}
   \end{minipage}}  & \texttt{result = get\_file(env)} \\
   &   & \texttt{expected = get\_file(cloud)} \\
   &   &  \texttt{rules = [\{"type":"sheet\_name"\},}\\
   &   &  \texttt{\ \  \  \  \  \  \  \ \ \{"type":"sheet\_data",} \\
   &   &  \texttt{\ \  \  \  \  \  \  \ \ \ "sheet\_idx0":0, } \\
   &   &  \texttt{\ \  \  \  \  \  \  \ \ \ "sheet\_idx1":1\}...]} \\
   &   &  \texttt{compare\_table(result, expected, rules)} \\
   &   & \\
    \midrule
    \multirow{4}{*}{\raisebox{-1cm}{\includegraphics[width=5cm]{images/evaluation_scripts_showcase/thunderbird_piad_tution.png}}} &  \multirow{4}{*}{
   \begin{minipage}{.35\textwidth}
   \textit{I've drafted an e-mail reminder for those who haven't paid tuition. Please help me to check out their e-mails from the payment record and add to the receiver field.}
   \end{minipage}} & \texttt{tree = get\_a11y\_tree(env)} \\
    &   & \texttt{rules = [\{"selectors": } \\
    &   & \texttt{\ \ ["tool-bar[attr|id=MsgHeadersToolbar]}\\ 
    &   & \texttt{\ \ label[name=To]~}\\
    &   & \texttt{\ \ [attr|class=\textbackslash"address-pill\textbackslash"]>}\\
    &   & \texttt{\ \ label[attr|class=\textbackslash"pill-label\textbackslash"]} \\
    &   & \texttt{\ \ [name*=\textbackslash"fox@someuniversity.edu...]} \\
    &   & \texttt{check\_a11y\_tree(tree, rules)} \\
  \bottomrule
\end{tabular}
}
\vspace{-18pt}
\label{tab:evaluation_examples}
\end{table}

\subsection{Observation Space}
\label{observation_space}
The observation space in \ours contains a \textbf{complete screenshot of the desktop screen}, including the mouse's position and shape, various application windows, files, and folders that are opened in different sizes and orders, maintaining the same perception as a human. 
Also, to be aligned with previous agent-building web and mobile research~\citep{liu2018reinforcement, li2020mapping, deng2023mind2web, zhou2023webarena} that provide and support the use of the webpage's DOM and app's view hierarchy, \ours also provides \textbf{XML-format accessibility (a11y) tree} (obtained via ATSPI~\footnote{\url{https://docs.gtk.org/atspi2/}} on Ubuntu, via PyWinAuto on Windows, \textit{etc.}), which can support additional information for modeling.
These raw observations allow rich interactions between multiple applications but induce challenges in long-horizon decision-making from high-resolution images (\textit{e.g.},~4k screenshots) and structured long text (\textit{e.g.},~accessibility trees).
For more detailed information on observation space, refer to App.~\ref{app:observation_space}.


\subsection{Action Space}
\label{action_space}

\begin{wraptable}{r}{8cm}
\vspace{-5pt}
\caption{Some examples of the mouse and keyboard actions $\mathcal{A}$ in \ours. 
See App.~\ref{app:action_space} for the complete list.}  
\centering
\resizebox{1.0\linewidth}{!}{%
\begin{tabular}{@{}ll@{}}
\toprule
Function & Description \\
\midrule
\texttt{moveTo(x, y)} & Moves the mouse to the specified coordinates. \\
\texttt{click(x, y)} & Clicks at the specified coordinates. \\
\texttt{write(`text')} & Types the specified text at the current cursor location. \\
\texttt{press(`enter')} & Presses the Enter key. \\
\texttt{hotkey(`ctrl', `c')} & Performs the Ctrl+C hotkey combination (copy). \\
\texttt{scroll(200)} & Scrolls up by 200 units. \\
\texttt{scroll(-200)} & Scrolls down by 200 units. \\
\texttt{dragTo(x, y)} & Drags the mouse to the specified coordinates. \\
\texttt{keyDown(`shift')} & Holds down the Shift key. \\
\texttt{keyUp(`shift')} & Releases the Shift key. \\
\texttt{WAIT} & Agent decides it should wait. \\
\texttt{FAIL} & Agent decides the task is infeasible. \\
\texttt{DONE} & Agent decides the task is finished. \\
\bottomrule
\end{tabular}
}
\vspace{-0.1in}
\label{tab:pyautogui_examples}
\end{wraptable}

Action space 
$\mathcal{A}$ in \ours encompasses all mouse and keyboard actions, including movement, clicks (left-key, right-key, multiple clicks), dragging, keystrokes, hotkeys, and others, covering all human-computer action space.
Some action examples are shown in Tab.~\ref{tab:pyautogui_examples} and the complete action list can be found in Appendix~\ref{app:action_space}. 
We use the widely used mouse and keyboard control library \texttt{pyautogui}\footnote{\url{https://pyautogui.readthedocs.io/en/latest/}} for our action space. 
This library leverages the high-level programming language Python to replicate and replay various human inputs into computers through code, allowing us to construct a universal and complete representation of actions. 
The agent must generate syntax-correct \texttt{pyautogui} Python code to predict valid actions. 
Basic actions, such as \texttt{press} and \texttt{moveTo}, can be integrated within program structures, such as for-loops, significantly improving the expressiveness of an action.
Timing is also crucial, as highlighted in previous studies on mobile devices~\citep{toyama2021androidenv}, as well as the ability to determine whether a task is infeasible or completed. 
Therefore, we add three special actions named \texttt{WAIT}, \texttt{FAIL}, and \texttt{DONE} to enhance the aforementioned action spaces.
Previous efforts towards creating domain-specific agents, such as MiniWoB++~\citep{shi2017world, liu2018reinforcement}, CC-Net~\citep{humphreys2022data}, and WebArena~\citep{zhou2023webarena, Koh2024VisualWebArenaEM}, have defined action spaces that include clicks and typing, as well as some actions specially designed for web browsing.
However, they do not model all possible actions on a computer, leading to limitations when attempting actions like right-clicking and clicking with the \texttt{ctrl} key held to select items.
This imposes an upper bound on agent learning capabilities.
